\chapter{Literature Review}\label{ch:literature}

\section{Modeling levels for particle-laden flow}
The hydrodynamics of a particle suspension can be described at different
levels of resolution. A particle-resolved method imposes boundary conditions
on each particle surface. It can represent the disturbance flow, finite-size
effects, and surface stresses when these are sufficiently resolved. Direct
forcing immersed-boundary methods provide one established route to such
calculations \citep{uhlmann2005}. Their additional resolution requirements
are justified when particle-scale flow is itself an essential part of the
problem.

A point-particle model instead advances particle centers using a force law
obtained under specified physical assumptions. The classical equation of
motion for a small rigid sphere includes effects beyond steady drag
\citep{maxey1983}. The heavy-particle, low-slip-Reynolds-number approximation
used here retains a much narrower description. It is appropriate to assess
this approximation as a simplified model rather than attribute to it the
full physics of the parent equation.

The strength of a point-particle approach is its economy of representation.
Its limits follow from the same reduction. A particle center and diameter
do not impose no-slip conditions on an actual moving surface. Finite-size
wall stresses, rotational dynamics, and lubrication do not appear unless
additional closures are included. These distinctions are particularly
important when a channel becomes narrow relative to particle diameter
\citep{clift1978}.

\section{One-way and two-way coupling}
In one-way coupling, the fluid field determines the particle force but the
particle reaction is neglected in the fluid equations. This separation is
useful for studying dilute transport because the fluid solution can be
evaluated independently of particle initial conditions. Small particle
volume fraction is a necessary consideration, although it does not alone
guarantee negligible feedback when the particle density is high. Mass
loading and the relevant time scales also matter.

Two-way coupling includes the hydrodynamic reaction in the carrier-fluid
momentum balance. The coupling must specify how a particle samples the fluid
and how its reaction is distributed to the Eulerian field. Treating either
operation independently can generate an inconsistent power exchange even
when the integrated force is correct. A shared transfer operator provides a
direct means to avoid this particular inconsistency.

Regularized point-particle methods provide a principled context for this
transfer. In the exact regularized point-particle method, the treatment of
the disturbance and its regularization is tied to the governing flow
equations \citep{gualtieri2015}. A generic normalized Gaussian deposition
does not inherit all of that method's properties. In particular, sampling
the disturbed velocity at a force source is not equivalent to recovering
the undisturbed velocity required by an isolated-particle drag law.

\section{Gaussian parcels and represented populations}
A parcel represents a population of physical particles through a
multiplicity. Its represented mass must scale with that multiplicity, and
its reaction force must use the same convention. A Gaussian kernel spreads
that force over a finite region, avoiding a singular source in the carrier
equations. The width of the kernel is a modeling parameter with a direct
effect on the spatial distribution of forcing.

Three lengths must be kept distinct: the physical diameter, the numerical
or model regularization width, and the local channel height. Equating the
Gaussian width with a particle radius without justification would blur
point-particle and finite-size interpretations. Likewise, changing the
width while keeping quadrature resolution fixed does not constitute an
ordinary grid-convergence study: it changes the regularized model being
solved.

Near a wall, a free-space Gaussian is truncated by the fluid domain. A
discrete normalization can preserve its integrated weight on the chosen
quadrature. That normalization establishes a transfer identity; it does
not derive a wall-corrected drag law. This distinction carries directly into
the constricted-channel extension developed later.

\section{Physics-informed representations}
Physics-informed neural networks approximate solution fields with
differentiable functions and train those functions using equation and
boundary residuals \citep{raissi2019}. They provide continuous field values
without a separate spatial interpolant. Automatic differentiation also
allows the same representation to supply the derivatives entering the
equations. These features are useful for particle transport, where the
evaluation locations move continuously.

Physics information in the objective is not a guarantee that optimization
finds a sufficiently accurate solution. Residual terms can have different
scales and different optimization behavior. Gradient imbalance and training
difficulties are documented limitations of PINNs
\citep{wang2021,wang2022}. A weighted objective may decrease while a
transport-sensitive region remains inaccurate. Consequently, reporting only
the optimization history would provide insufficient support for the
predictions considered here.

Hard constraints offer an additional design choice. An ansatz can impose an
initial condition, periodicity, or homogeneous wall velocity exactly, leaving
the network to represent a correction that respects those conditions. This
reduces the number of requirements enforced through a penalty. It does not
remove the need to evaluate momentum, continuity, and particle-motion
residuals at independent points.

The baseline and Gaussian stages use different neural roles. The baseline
flow network approximates a steady velocity--pressure field, whereas the
force network approximates a prescribed algebraic law. In the Gaussian
stage, separate fluid and parcel networks are coupled through common
residuals. This is a continuous-time coupled forward representation rather
than a trajectory regression against reference data.

\section{Independent references and numerical credibility}
The finite-element method supplies an independently discretized reference
for the baseline carrier flow. Taylor--Hood velocity--pressure elements are
a standard mixed formulation for incompressible flow \citep{taylor1973}.
Their assembly here uses \texttt{scikit-fem} \citep{gustafsson2020}. An
independent discretization is valuable because its numerical errors differ
from those of the neural representation. It must still be evaluated for
nonlinear convergence, incompressibility, and mesh sensitivity.

Verification and physical validation address different questions
\citep{roache1998}. Agreement with an analytical solution tests whether a
specified mathematical problem is solved correctly. Agreement with another
numerical solver tests consistency for a matched model and boundary-value
problem. Experimental validation additionally requires suitable
measurements, controlled conditions, and compatible physical assumptions.
The results in this thesis primarily concern the first two questions.

A relative velocity error should be interpreted together with local
discrepancies, boundary errors, flux variation, and particle-motion
diagnostics. Pressure requires particular care because incompressible
velocity boundary conditions determine pressure only up to an additive
constant. Comparing unaligned pressure arrays can produce a large error
without a corresponding physical disagreement. Conversely, an error
measured on a section where the reference pressure is constant can have an
ill-conditioned denominator. A spatially extended, gauge-aligned comparison
avoids these problems.

\section{Transport accuracy and confinement}
Particle endpoints depend on a time-integrated response to the carrier
field. A small discrepancy may accumulate along a long path or determine
whether a particle approaches a capture boundary. Near-wall outcomes can
therefore be less accurate than bulk flow metrics suggest. A meaningful
trajectory comparison must pair the same initial conditions, restrict its
position error to the common time interval, and separately report whether
the two solutions reach the same endpoint class.

Confinement adds a physical limitation to this numerical sensitivity. When
the local channel height approaches the particle diameter, the unbounded
Stokes law becomes a questionable description even if the numerical
trajectory is well resolved. A study of narrowing should therefore track
both numerical agreement and the changing ratio of particle diameter to
local gap. A dilute regularized model cannot be used to infer dense
blockage merely because a narrow geometry is drawn.

\section{Position of the present study}
\Cref{tab:model_comparison} summarizes the modeling levels relevant to the
study. The purpose is to build a verified chain across the first two levels
and specify the requirements for a constrained extension. The focus is
neither a new empirical drag correlation nor a replacement for resolved
particle hydrodynamics. The main methodological question is whether a
neural representation can support the stated fluid and particle quantities
while preserving a consistent transfer of force and work.

\begin{table}[htbp]\centering\small
\caption{Model classes and the evidence they require.}\label{tab:model_comparison}
\begin{tabularx}{\textwidth}{@{}p{.22\textwidth}YY@{}}\toprule
Model & Main representation & Essential assessment\\\midrule
One-way point particles & Carrier field and particle ODEs & Field error, slip regime, trajectory agreement\\
Gaussian parcels & Population-weighted regularized forcing & Shared transfer, residuals, budgets, width sensitivity\\
Resolved particles & Moving solid surfaces and surrounding flow & Surface resolution, forces, torques and gap flow\\
Constricted dilute model & Variable-height fluid domain with particles & Wall geometry, flux, throat resolution and closure validity\\
\bottomrule\end{tabularx}
\end{table}
